\chapter{Arquivos do projeto}
\label{ap:arquivos}

O trabalho foi organizado em quatro projetos independentes do Quartus II, um por pasta de \texttt{project/}, cada uma com a mesma estrutura (substituir \texttt{<nome>} por \texttt{reg\_bank}, \texttt{reg\_status}, \texttt{stack} ou \texttt{prog\_cnt}):

\begin{description}
  \item[\texttt{<nome>.vhd}] descrição VHDL do bloco, apresentada nos Códigos~\ref{lst:reg-bank}, \ref{lst:reg-status}, \ref{lst:stack} e \ref{lst:prog-cnt}. O relatório aponta direto para esses arquivos.
  \item[\texttt{<nome>.qpf}, \texttt{<nome>.qsf}] arquivos do projeto no Quartus II, com o dispositivo (Cyclone II EP2C35F672C6), o arquivo VHDL e todos os arquivos de forma de onda.
  \item[\texttt{<nome>.sdc}] restrição de clock de 50\,MHz usada pelo TimeQuest.
  \item[\texttt{<nome>*.vwf}] formas de onda de simulação; \texttt{<nome>.vwf} é a visão geral e os demais, os casos específicos, descritos nos capítulos de cada bloco.
  \item[\texttt{<nome>\_simulation.png}] captura da simulação de \texttt{<nome>.vwf}, usada nas Figuras~\ref{fig:sim-reg-bank}, \ref{fig:sim-reg-status}, \ref{fig:sim-stack} e \ref{fig:sim-prog-cnt}.
  \item[\texttt{<nome>\_rtl.svg}] circuito em nível RTL gerado pelo RapidRTL, usado nas Figuras~\ref{fig:rtl-reg-bank}, \ref{fig:rtl-reg-status}, \ref{fig:rtl-stack} e \ref{fig:rtl-prog-cnt}. O \texttt{justfile} do relatório converte cada um em PDF com o Inkscape.
  \item[\texttt{extras/<nome>\_tb.vhd}] \textit{testbench} auto-verificável, executado com o \texttt{nvc}.
  \item[\texttt{extras/casos\_de\_teste.md}] tabelas de resultados esperados de todas as formas de onda do bloco.
\end{description}

Os \textit{testbenches} são executados, dentro da pasta do bloco, com:

\begingroup\footnotesize
\begin{verbatim}
nvc --std=1993 -H 512m -a <nome>.vhd extras/<nome>_tb.vhd -e <nome>_tb -r
\end{verbatim}
\endgroup

O relatório fica em \texttt{report/tp2/} e é gerado com \texttt{just pdf}. Os blocos de código usam o pacote \texttt{minted}, e por isso a compilação exige o Pygments e a opção \texttt{-shell-escape}, já incluída no \texttt{justfile}.

\chapter{Casos de teste das demais formas de onda}
\label{ap:tabelas}

Este apêndice traz os valores esperados das formas de onda que não aparecem nos capítulos de cada bloco. Cada linha é um intervalo entre mudanças das entradas, e as saídas são as do fim do intervalo, depois de qualquer borda de subida dentro dele. As formas de onda de \texttt{reg\_bank}, \texttt{reg\_status}, \texttt{stack} e \texttt{prog\_cnt} são apresentadas nessa ordem.

\section{Banco de registradores}
{\scriptsize\setlength{\tabcolsep}{3pt}
\begin{longtable}{rcccccccp{5.3cm}}
\caption{Casos de teste de \texttt{reg\_bank\_escrita.vwf}: escrita nos quatro registradores.}
\label{tab:caso-reg-bank-escrita}\\
\hline $t$ (ns) & \texttt{nrst} & \texttt{di} & \texttt{a\_sel} & \texttt{b\_sel} & \texttt{wr\_ena} & \texttt{a\_do} & \texttt{b\_do} & Observação \\ \hline \endfirsthead
\hline $t$ (ns) & \texttt{nrst} & \texttt{di} & \texttt{a\_sel} & \texttt{b\_sel} & \texttt{wr\_ena} & \texttt{a\_do} & \texttt{b\_do} & Observação \\ \hline \endhead
0--20 & 0 & 00h & 0 & 0 & 0 & 00h & 00h &  \\
20--40 & 1 & A5h & 0 & 0 & 1 & A5h & A5h & escreve A5h em R0 (b\_sel = a\_sel) \\
40--60 & 1 & 5Ah & 1 & 1 & 1 & 5Ah & 5Ah & escreve 5Ah em R1 \\
60--80 & 1 & FFh & 2 & 2 & 1 & FFh & FFh & escreve FFh em R2 \\
80--100 & 1 & 01h & 3 & 3 & 1 & 01h & 01h & escreve 01h em R3 \\
100--120 & 1 & 00h & 0 & 3 & 0 & A5h & 01h & desabilita a escrita; a\_do = R0, b\_do = R3 \\
120--140 & 1 & 00h & 1 & 2 & 0 & 5Ah & FFh & a\_do = R1, b\_do = R2 \\
140--160 & 1 & 00h & 2 & 1 & 0 & FFh & 5Ah & a\_do = R2, b\_do = R1 \\
160--180 & 1 & 00h & 3 & 0 & 0 & 01h & A5h & a\_do = R3, b\_do = R0 \\
180--200 & 1 & 00h & 0 & 0 & 1 & 00h & 00h & sobrescreve R0 com 00h \\
200--220 & 1 & FFh & 3 & 3 & 1 & FFh & FFh & sobrescreve R3 com FFh \\
220--240 & 1 & FFh & 0 & 3 & 0 & 00h & FFh & lê R0 e R3 \\
240--260 & 1 & FFh & 1 & 2 & 0 & 5Ah & FFh & a\_do = R1, b\_do = R2 (intactos) \\
\hline
\end{longtable}}
\vspace{-0.9cm}
\fonte{Autoria própria. Saídas do fim de cada intervalo, conforme o modelo de referência e a simulação do Quartus II.}

{\scriptsize\setlength{\tabcolsep}{3pt}
\begin{longtable}{rcccccccp{5.3cm}}
\caption{Casos de teste de \texttt{reg\_bank\_leitura.vwf}: leitura assíncrona.}
\label{tab:caso-reg-bank-leitura}\\
\hline $t$ (ns) & \texttt{nrst} & \texttt{di} & \texttt{a\_sel} & \texttt{b\_sel} & \texttt{wr\_ena} & \texttt{a\_do} & \texttt{b\_do} & Observação \\ \hline \endfirsthead
\hline $t$ (ns) & \texttt{nrst} & \texttt{di} & \texttt{a\_sel} & \texttt{b\_sel} & \texttt{wr\_ena} & \texttt{a\_do} & \texttt{b\_do} & Observação \\ \hline \endhead
0--20 & 0 & 00h & 0 & 0 & 0 & 00h & 00h &  \\
20--40 & 1 & 10h & 0 & 0 & 1 & 10h & 10h & R0 = 10h \\
40--60 & 1 & 20h & 1 & 0 & 1 & 20h & 10h & R1 = 20h \\
60--80 & 1 & 40h & 2 & 0 & 1 & 40h & 10h & R2 = 40h \\
80--100 & 1 & 80h & 3 & 0 & 1 & 80h & 10h & R3 = 80h \\
100--105 & 1 & FFh & 3 & 0 & 0 & 80h & 10h & wr\_ena = 0, di = FFh (ignorado) \\
105--125 & 1 & FFh & 0 & 0 & 0 & 10h & 10h & a\_sel = 0, b\_sel = 0 \\
125--145 & 1 & FFh & 0 & 1 & 0 & 10h & 20h & a\_sel = 0, b\_sel = 1 \\
145--165 & 1 & FFh & 0 & 2 & 0 & 10h & 40h & a\_sel = 0, b\_sel = 2 \\
165--185 & 1 & FFh & 0 & 3 & 0 & 10h & 80h & a\_sel = 0, b\_sel = 3 \\
185--205 & 1 & FFh & 1 & 0 & 0 & 20h & 10h & a\_sel = 1, b\_sel = 0 \\
205--225 & 1 & FFh & 1 & 1 & 0 & 20h & 20h & a\_sel = 1, b\_sel = 1 \\
225--245 & 1 & FFh & 1 & 2 & 0 & 20h & 40h & a\_sel = 1, b\_sel = 2 \\
245--265 & 1 & FFh & 1 & 3 & 0 & 20h & 80h & a\_sel = 1, b\_sel = 3 \\
265--285 & 1 & FFh & 2 & 0 & 0 & 40h & 10h & a\_sel = 2, b\_sel = 0 \\
285--305 & 1 & FFh & 2 & 1 & 0 & 40h & 20h & a\_sel = 2, b\_sel = 1 \\
305--325 & 1 & FFh & 2 & 2 & 0 & 40h & 40h & a\_sel = 2, b\_sel = 2 \\
325--345 & 1 & FFh & 2 & 3 & 0 & 40h & 80h & a\_sel = 2, b\_sel = 3 \\
345--365 & 1 & FFh & 3 & 0 & 0 & 80h & 10h & a\_sel = 3, b\_sel = 0 \\
365--385 & 1 & FFh & 3 & 1 & 0 & 80h & 20h & a\_sel = 3, b\_sel = 1 \\
385--405 & 1 & FFh & 3 & 2 & 0 & 80h & 40h & a\_sel = 3, b\_sel = 2 \\
405--420 & 1 & FFh & 3 & 3 & 0 & 80h & 80h & a\_sel = 3, b\_sel = 3 \\
\hline
\end{longtable}}
\vspace{-0.9cm}
\fonte{Autoria própria. Saídas do fim de cada intervalo, conforme o modelo de referência e a simulação do Quartus II.}

{\scriptsize\setlength{\tabcolsep}{3pt}
\begin{longtable}{rcccccccp{5.3cm}}
\caption{Casos de teste de \texttt{reg\_bank\_wr\_ena.vwf}: habilitação da escrita.}
\label{tab:caso-reg-bank-wr-ena}\\
\hline $t$ (ns) & \texttt{nrst} & \texttt{di} & \texttt{a\_sel} & \texttt{b\_sel} & \texttt{wr\_ena} & \texttt{a\_do} & \texttt{b\_do} & Observação \\ \hline \endfirsthead
\hline $t$ (ns) & \texttt{nrst} & \texttt{di} & \texttt{a\_sel} & \texttt{b\_sel} & \texttt{wr\_ena} & \texttt{a\_do} & \texttt{b\_do} & Observação \\ \hline \endhead
0--20 & 0 & 00h & 0 & 0 & 0 & 00h & 00h &  \\
20--40 & 1 & 11h & 0 & 0 & 1 & 11h & 11h & escreve 11h em R0 \\
40--60 & 1 & 99h & 0 & 0 & 0 & 11h & 11h & wr\_ena = 0, di = 99h: R0 não muda \\
60--80 & 1 & 88h & 0 & 0 & 0 & 11h & 11h & di = 88h: R0 continua 11h \\
80--100 & 1 & 99h & 0 & 0 & 1 & 99h & 99h & wr\_ena = 1: escreve 99h em R0 \\
100--120 & 1 & 00h & 0 & 0 & 0 & 99h & 99h & wr\_ena = 0 \\
120--132 & 1 & 66h & 1 & 1 & 0 & 00h & 00h & prepara R1 = 66h, mas sem escrever \\
132--138 & 1 & 66h & 1 & 1 & 1 & 00h & 00h & pulso de wr\_ena entre duas bordas (12 a 18 ns): não escreve \\
138--148 & 1 & 66h & 1 & 1 & 0 & 00h & 00h & fim do pulso \\
148--152 & 1 & 66h & 1 & 1 & 1 & 66h & 66h & pulso de wr\_ena sobre a borda de subida (8 a 12 ns): escreve 66h em R1 \\
152--160 & 1 & 66h & 1 & 1 & 0 & 66h & 66h & fim do pulso \\
160--180 & 1 & 66h & 1 & 0 & 0 & 66h & 99h & lê R1 e R0 \\
180--200 & 1 & 66h & 0 & 1 & 0 & 99h & 66h & lê R0 e R1 \\
\hline
\end{longtable}}
\vspace{-0.9cm}
\fonte{Autoria própria. Saídas do fim de cada intervalo, conforme o modelo de referência e a simulação do Quartus II.}

{\scriptsize\setlength{\tabcolsep}{3pt}
\begin{longtable}{rcccccccp{5.3cm}}
\caption{Casos de teste de \texttt{reg\_bank\_reset.vwf}: reset assíncrono.}
\label{tab:caso-reg-bank-reset}\\
\hline $t$ (ns) & \texttt{nrst} & \texttt{di} & \texttt{a\_sel} & \texttt{b\_sel} & \texttt{wr\_ena} & \texttt{a\_do} & \texttt{b\_do} & Observação \\ \hline \endfirsthead
\hline $t$ (ns) & \texttt{nrst} & \texttt{di} & \texttt{a\_sel} & \texttt{b\_sel} & \texttt{wr\_ena} & \texttt{a\_do} & \texttt{b\_do} & Observação \\ \hline \endhead
0--20 & 0 & 00h & 0 & 0 & 0 & 00h & 00h &  \\
20--40 & 1 & C1h & 0 & 0 & 1 & C1h & C1h & escreve C1h em R0 \\
40--60 & 1 & C2h & 1 & 0 & 1 & C2h & C1h & escreve C2h em R1 \\
60--80 & 1 & C3h & 2 & 0 & 1 & C3h & C1h & escreve C3h em R2 \\
80--85 & 1 & C3h & 0 & 1 & 0 & C1h & C2h & a\_do = R0, b\_do = R1; wr\_ena = 0 \\
85--100 & 0 & C3h & 0 & 1 & 0 & 00h & 00h & reset assíncrono (fora de qualquer borda): tudo vai a 0 \\
100--123 & 0 & EEh & 1 & 1 & 1 & 00h & 00h & ainda em reset; wr\_ena = 1 sobre a borda: não escreve \\
123--140 & 1 & EEh & 1 & 1 & 1 & EEh & EEh & solta o reset (a borda em 130 ns já pode escrever) \\
140--160 & 1 & EEh & 0 & 1 & 0 & 00h & EEh & wr\_ena = 0 \\
160--180 & 1 & D1h & 0 & 1 & 1 & D1h & EEh & escreve D1h em R0 \\
180--212 & 1 & D1h & 0 & 1 & 0 & D1h & EEh & wr\_ena = 0 \\
212--215 & 0 & D1h & 0 & 1 & 0 & 00h & 00h & pulso de reset de 3 ns, sem borda: ainda zera \\
215--220 & 1 & D1h & 0 & 1 & 0 & 00h & 00h & fim do pulso \\
220--240 & 1 & D1h & 0 & 1 & 0 & 00h & 00h & tudo em 0 \\
\hline
\end{longtable}}
\vspace{-0.9cm}
\fonte{Autoria própria. Saídas do fim de cada intervalo, conforme o modelo de referência e a simulação do Quartus II.}


\section{Registrador de status}
{\scriptsize\setlength{\tabcolsep}{3pt}
\begin{longtable}{rccccp{4.4cm}}
\caption{Casos de teste de \texttt{reg\_status\_escrita.vwf}: escrita de todos os bits.}
\label{tab:caso-reg-status-escrita}\\
\hline $t$ (ns) & \texttt{nrst} & \texttt{status\_reg\_in} & \texttt{status\_reg\_wr\_ena} & \texttt{status\_reg\_out} & Observação \\ \hline \endfirsthead
\hline $t$ (ns) & \texttt{nrst} & \texttt{status\_reg\_in} & \texttt{status\_reg\_wr\_ena} & \texttt{status\_reg\_out} & Observação \\ \hline \endhead
0--20 & 0 & 00000 & 00000 & 00000 &  \\
20--40 & 1 & 10101 & 11111 & 10101 & escreve 10101 \\
40--60 & 1 & 01010 & 11111 & 01010 & escreve 01010 \\
60--80 & 1 & 11111 & 11111 & 11111 & escreve 11111 \\
80--100 & 1 & 00000 & 11111 & 00000 & escreve 00000 \\
100--120 & 1 & 11111 & 11111 & 11111 & escreve 11111 \\
120--140 & 1 & 00110 & 11111 & 00110 & escreve 00110 \\
140--160 & 1 & 11001 & 00000 & 00110 & mantém 00110 (clock continua) \\
\hline
\end{longtable}}
\vspace{-0.9cm}
\fonte{Autoria própria. Saídas do fim de cada intervalo, conforme o modelo de referência e a simulação do Quartus II.}

{\scriptsize\setlength{\tabcolsep}{3pt}
\begin{longtable}{rccccp{4.4cm}}
\caption{Casos de teste de \texttt{reg\_status\_enable.vwf}: habilitação bit a bit.}
\label{tab:caso-reg-status-enable}\\
\hline $t$ (ns) & \texttt{nrst} & \texttt{status\_reg\_in} & \texttt{status\_reg\_wr\_ena} & \texttt{status\_reg\_out} & Observação \\ \hline \endfirsthead
\hline $t$ (ns) & \texttt{nrst} & \texttt{status\_reg\_in} & \texttt{status\_reg\_wr\_ena} & \texttt{status\_reg\_out} & Observação \\ \hline \endhead
0--20 & 0 & 00000 & 00000 & 00000 &  \\
20--40 & 1 & 11111 & 00001 & 00001 & seta só RN \\
40--60 & 1 & 11111 & 00010 & 00011 & seta só RV \\
60--80 & 1 & 11111 & 00100 & 00111 & seta só RZ \\
80--100 & 1 & 11111 & 01000 & 01111 & seta só RH \\
100--120 & 1 & 11111 & 10000 & 11111 & seta só RC \\
120--140 & 1 & 00000 & 00000 & 11111 & ninguém habilitado, in = 00000: mantém 11111 \\
140--160 & 1 & 00000 & 00001 & 11110 & zera só RN \\
160--180 & 1 & 00000 & 00010 & 11100 & zera só RV \\
180--200 & 1 & 00000 & 00100 & 11000 & zera só RZ \\
200--220 & 1 & 00000 & 01000 & 10000 & zera só RH \\
220--240 & 1 & 00000 & 10000 & 00000 & zera só RC \\
240--260 & 1 & 10101 & 10101 & 10101 & habilita RN, RZ, RC com in = 10101 \\
260--280 & 1 & 01010 & 11010 & 01111 & habilita os outros dois com in = 01010 e RC com 0 \\
\hline
\end{longtable}}
\vspace{-0.9cm}
\fonte{Autoria própria. Saídas do fim de cada intervalo, conforme o modelo de referência e a simulação do Quartus II.}

{\scriptsize\setlength{\tabcolsep}{3pt}
\begin{longtable}{rccccp{4.4cm}}
\caption{Casos de teste de \texttt{reg\_status\_alu.vwf}: uso com a ALU (flags C H Z V N).}
\label{tab:caso-reg-status-alu}\\
\hline $t$ (ns) & \texttt{nrst} & \texttt{status\_reg\_in} & \texttt{status\_reg\_wr\_ena} & \texttt{status\_reg\_out} & Observação \\ \hline \endfirsthead
\hline $t$ (ns) & \texttt{nrst} & \texttt{status\_reg\_in} & \texttt{status\_reg\_wr\_ena} & \texttt{status\_reg\_out} & Observação \\ \hline \endhead
0--20 & 0 & 00000 & 00000 & 00000 &  \\
20--40 & 1 & 11100 & 11111 & 11100 & ADD 77h+89h = 00h: C=1 H=1 Z=1 V=0 N=0, todos habilitados \\
40--60 & 1 & 00000 & 00101 & 11000 & AND resultado 30h: só N=0 e Z=0 são atualizados (C, H, V preservados) \\
60--80 & 1 & 01010 & 11111 & 01010 & SUB 80h-01h = 7Fh: C=0 H=1 Z=0 V=1 N=0, todos habilitados \\
80--100 & 1 & 00100 & 00101 & 01110 & XOR resultado 00h: só N=0 e Z=1 \\
100--120 & 1 & 10001 & 10101 & 11011 & SRA 81h = C0h: C=1 Z=0 N=1 (H e V preservados) \\
120--140 & 1 & 11111 & 00000 & 11011 & nenhum flag atualizado (instrução que não altera flags) \\
140--160 & 1 & 00000 & 11111 & 00000 & ADD 01h+01h = 02h: tudo 0 \\
160--180 & 1 & 10001 & 11111 & 10001 & ADD C7h+D0h = 97h: C=1 H=0 Z=0 V=0 N=1 \\
180--200 & 1 & 10001 & 11111 & 10001 & mantém \\
\hline
\end{longtable}}
\vspace{-0.9cm}
\fonte{Autoria própria. Saídas do fim de cada intervalo, conforme o modelo de referência e a simulação do Quartus II.}

{\scriptsize\setlength{\tabcolsep}{3pt}
\begin{longtable}{rccccp{4.4cm}}
\caption{Casos de teste de \texttt{reg\_status\_reset.vwf}: reset assíncrono.}
\label{tab:caso-reg-status-reset}\\
\hline $t$ (ns) & \texttt{nrst} & \texttt{status\_reg\_in} & \texttt{status\_reg\_wr\_ena} & \texttt{status\_reg\_out} & Observação \\ \hline \endfirsthead
\hline $t$ (ns) & \texttt{nrst} & \texttt{status\_reg\_in} & \texttt{status\_reg\_wr\_ena} & \texttt{status\_reg\_out} & Observação \\ \hline \endhead
0--20 & 0 & 00000 & 00000 & 00000 &  \\
20--40 & 1 & 11111 & 11111 & 11111 & escreve 11111 \\
40--65 & 1 & 11111 & 00000 & 11111 & mantém \\
65--80 & 0 & 11111 & 00000 & 00000 & reset assíncrono, fora de qualquer borda \\
80--103 & 0 & 11111 & 11111 & 00000 & ainda em reset; escrita habilitada sobre a borda: não escreve \\
103--120 & 1 & 11111 & 11111 & 11111 & solta o reset \\
120--140 & 1 & 10101 & 11111 & 10101 & escreve 10101 (a borda desse ciclo já escreve) \\
140--172 & 1 & 10101 & 00000 & 10101 & mantém \\
172--175 & 0 & 10101 & 00000 & 00000 & pulso de reset de 3 ns, sem borda \\
175--180 & 1 & 10101 & 00000 & 00000 & fim do pulso \\
180--200 & 1 & 10101 & 00000 & 00000 & tudo em 0 \\
\hline
\end{longtable}}
\vspace{-0.9cm}
\fonte{Autoria própria. Saídas do fim de cada intervalo, conforme o modelo de referência e a simulação do Quartus II.}


\section{Pilha}
{\scriptsize\setlength{\tabcolsep}{3pt}
\begin{longtable}{rcccccp{5.7cm}}
\caption{Casos de teste de \texttt{stack\_lifo.vwf}: ordem LIFO.}
\label{tab:caso-stack-lifo}\\
\hline $t$ (ns) & \texttt{nrst} & \texttt{stack\_in} & \texttt{stack\_push} & \texttt{stack\_pop} & \texttt{stack\_out} & Observação \\ \hline \endfirsthead
\hline $t$ (ns) & \texttt{nrst} & \texttt{stack\_in} & \texttt{stack\_push} & \texttt{stack\_pop} & \texttt{stack\_out} & Observação \\ \hline \endhead
0--20 & 0 & 0 & 0 & 0 & 0 &  \\
20--40 & 1 & 100 & 1 & 0 & 100 & push 100 \\
40--60 & 1 & 200 & 1 & 0 & 200 & push 200 \\
60--80 & 1 & 300 & 1 & 0 & 300 & push 300 \\
80--100 & 1 & 400 & 1 & 0 & 400 & push 400 \\
100--120 & 1 & 500 & 1 & 0 & 500 & push 500 \\
120--140 & 1 & 600 & 1 & 0 & 600 & push 600 \\
140--160 & 1 & 700 & 1 & 0 & 700 & push 700 \\
160--180 & 1 & 800 & 1 & 0 & 800 & push 800 \\
180--200 & 1 & 0 & 0 & 1 & 700 & pop \\
200--220 & 1 & 0 & 0 & 1 & 600 & pop \\
220--240 & 1 & 0 & 0 & 1 & 500 & pop \\
240--260 & 1 & 0 & 0 & 1 & 400 & pop \\
260--280 & 1 & 0 & 0 & 1 & 300 & pop \\
280--300 & 1 & 0 & 0 & 1 & 200 & pop \\
300--320 & 1 & 0 & 0 & 1 & 100 & pop \\
320--340 & 1 & 0 & 0 & 1 & 0 & pop \\
340--360 & 1 & 0 & 0 & 1 & 0 & pop extra: pilha vazia \\
360--380 & 1 & 0 & 0 & 0 & 0 & sem operação \\
\hline
\end{longtable}}
\vspace{-0.9cm}
\fonte{Autoria própria. Saídas do fim de cada intervalo, conforme o modelo de referência e a simulação do Quartus II.}

{\scriptsize\setlength{\tabcolsep}{3pt}
\begin{longtable}{rcccccp{5.7cm}}
\caption{Casos de teste de \texttt{stack\_overflow.vwf}: stack overflow.}
\label{tab:caso-stack-overflow}\\
\hline $t$ (ns) & \texttt{nrst} & \texttt{stack\_in} & \texttt{stack\_push} & \texttt{stack\_pop} & \texttt{stack\_out} & Observação \\ \hline \endfirsthead
\hline $t$ (ns) & \texttt{nrst} & \texttt{stack\_in} & \texttt{stack\_push} & \texttt{stack\_pop} & \texttt{stack\_out} & Observação \\ \hline \endhead
0--20 & 0 & 0 & 0 & 0 & 0 &  \\
20--40 & 1 & 1 & 1 & 0 & 1 & push 1 \\
40--60 & 1 & 2 & 1 & 0 & 2 & push 2 \\
60--80 & 1 & 3 & 1 & 0 & 3 & push 3 \\
80--100 & 1 & 4 & 1 & 0 & 4 & push 4 \\
100--120 & 1 & 5 & 1 & 0 & 5 & push 5 \\
120--140 & 1 & 6 & 1 & 0 & 6 & push 6 \\
140--160 & 1 & 7 & 1 & 0 & 7 & push 7 \\
160--180 & 1 & 8 & 1 & 0 & 8 & push 8 \\
180--200 & 1 & 9 & 1 & 0 & 9 & push 9 \\
200--220 & 1 & 10 & 1 & 0 & 10 & push 10 \\
220--240 & 1 & 0 & 0 & 1 & 9 & pop \\
240--260 & 1 & 0 & 0 & 1 & 8 & pop \\
260--280 & 1 & 0 & 0 & 1 & 7 & pop \\
280--300 & 1 & 0 & 0 & 1 & 6 & pop \\
300--320 & 1 & 0 & 0 & 1 & 5 & pop \\
320--340 & 1 & 0 & 0 & 1 & 4 & pop \\
340--360 & 1 & 0 & 0 & 1 & 3 & pop \\
360--380 & 1 & 0 & 0 & 1 & 0 & pop \\
380--400 & 1 & 0 & 0 & 1 & 0 & pop \\
400--420 & 1 & 0 & 0 & 1 & 0 & pop \\
420--440 & 1 & 0 & 0 & 0 & 0 & sem operação \\
\hline
\end{longtable}}
\vspace{-0.9cm}
\fonte{Autoria própria. Saídas do fim de cada intervalo, conforme o modelo de referência e a simulação do Quartus II.}

{\scriptsize\setlength{\tabcolsep}{3pt}
\begin{longtable}{rcccccp{5.7cm}}
\caption{Casos de teste de \texttt{stack\_push\_pop.vwf}: push e pop simultâneos.}
\label{tab:caso-stack-push-pop}\\
\hline $t$ (ns) & \texttt{nrst} & \texttt{stack\_in} & \texttt{stack\_push} & \texttt{stack\_pop} & \texttt{stack\_out} & Observação \\ \hline \endfirsthead
\hline $t$ (ns) & \texttt{nrst} & \texttt{stack\_in} & \texttt{stack\_push} & \texttt{stack\_pop} & \texttt{stack\_out} & Observação \\ \hline \endhead
0--20 & 0 & 0 & 0 & 0 & 0 &  \\
20--40 & 1 & 10 & 1 & 0 & 10 & push 10 \\
40--60 & 1 & 20 & 1 & 0 & 20 & push 20 \\
60--80 & 1 & 30 & 1 & 1 & 30 & push + pop com 30: empilha 30 \\
80--100 & 1 & 40 & 1 & 1 & 40 & push + pop com 40: empilha 40 \\
100--120 & 1 & 0 & 0 & 1 & 30 & só pop (volta a 30) \\
120--140 & 1 & 0 & 0 & 1 & 20 & só pop (volta a 20) \\
140--160 & 1 & 0 & 0 & 1 & 10 & só pop (volta a 10) \\
160--180 & 1 & 0 & 0 & 1 & 0 & só pop (esvazia) \\
180--200 & 1 & 55 & 1 & 1 & 55 & push + pop com 55 na pilha vazia: empilha \\
200--220 & 1 & 0 & 0 & 1 & 0 & só pop (esvazia) \\
220--232 & 1 & 66 & 0 & 0 & 0 & pulso de push sem borda (12 a 18 ns): ignorado \\
232--238 & 1 & 66 & 1 & 0 & 0 & sobe stack\_push \\
238--240 & 1 & 66 & 0 & 0 & 0 & desce stack\_push \\
240--248 & 1 & 66 & 0 & 0 & 0 & pulso de push sobre a borda (8 a 12 ns): empilha 66 \\
248--252 & 1 & 66 & 1 & 0 & 66 & sobe stack\_push \\
252--260 & 1 & 66 & 0 & 0 & 66 & desce stack\_push \\
260--280 & 1 & 66 & 0 & 0 & 66 & sem operação \\
\hline
\end{longtable}}
\vspace{-0.9cm}
\fonte{Autoria própria. Saídas do fim de cada intervalo, conforme o modelo de referência e a simulação do Quartus II.}

{\scriptsize\setlength{\tabcolsep}{3pt}
\begin{longtable}{rcccccp{5.7cm}}
\caption{Casos de teste de \texttt{stack\_reset.vwf}: reset assíncrono.}
\label{tab:caso-stack-reset}\\
\hline $t$ (ns) & \texttt{nrst} & \texttt{stack\_in} & \texttt{stack\_push} & \texttt{stack\_pop} & \texttt{stack\_out} & Observação \\ \hline \endfirsthead
\hline $t$ (ns) & \texttt{nrst} & \texttt{stack\_in} & \texttt{stack\_push} & \texttt{stack\_pop} & \texttt{stack\_out} & Observação \\ \hline \endhead
0--20 & 0 & 0 & 0 & 0 & 0 &  \\
20--40 & 1 & 11 & 1 & 0 & 11 & push 11 \\
40--60 & 1 & 22 & 1 & 0 & 22 & push 22 \\
60--80 & 1 & 33 & 1 & 0 & 33 & push 33 \\
80--85 & 1 & 33 & 0 & 0 & 33 & sem operação \\
85--100 & 0 & 33 & 0 & 0 & 0 & reset assíncrono, fora de qualquer borda \\
100--123 & 0 & 44 & 1 & 0 & 0 & ainda em reset; push sobre a borda: não empilha \\
123--140 & 1 & 44 & 1 & 0 & 44 & solta o reset \\
140--160 & 1 & 77 & 1 & 0 & 77 & push 77 (a borda já empilha) \\
160--180 & 1 & 77 & 0 & 1 & 44 & pop (o topo vira 0: nada sobrou do que veio antes do reset) \\
180--200 & 1 & 77 & 0 & 1 & 0 & pop \\
200--212 & 1 & 77 & 0 & 0 & 0 & sem operação \\
212--215 & 0 & 77 & 0 & 0 & 0 & pulso de reset de 3 ns, sem borda \\
215--220 & 1 & 77 & 0 & 0 & 0 & fim do pulso \\
220--240 & 1 & 5 & 1 & 0 & 5 & push 5 \\
240--260 & 1 & 5 & 0 & 0 & 5 & sem operação \\
\hline
\end{longtable}}
\vspace{-0.9cm}
\fonte{Autoria própria. Saídas do fim de cada intervalo, conforme o modelo de referência e a simulação do Quartus II.}


\section{Contador de programa}
{\scriptsize\setlength{\tabcolsep}{3pt}
\begin{longtable}{rccccccp{4.2cm}}
\caption{Casos de teste de \texttt{prog\_cnt\_incremento.vwf}: incremento e volta módulo 2048.}
\label{tab:caso-prog-cnt-incremento}\\
\hline $t$ (ns) & \texttt{nrst} & \texttt{pc\_ctrl} & \texttt{new\_pc\_in} & \texttt{from\_stack} & \texttt{next\_pc\_out} & \texttt{pc\_out} & Observação \\ \hline \endfirsthead
\hline $t$ (ns) & \texttt{nrst} & \texttt{pc\_ctrl} & \texttt{new\_pc\_in} & \texttt{from\_stack} & \texttt{next\_pc\_out} & \texttt{pc\_out} & Observação \\ \hline \endhead
0--20 & 0 & 00 & 0 & 0 & 0 & 0 &  \\
20--40 & 1 & 10 & 2045 & 0 & 2045 & 2045 & carrega 2045 \\
40--60 & 1 & 11 & 2045 & 0 & 2047 & 2046 & incrementa \\
60--80 & 1 & 11 & 2045 & 0 & 0 & 2047 & incrementa \\
80--100 & 1 & 11 & 2045 & 0 & 1 & 0 & incrementa \\
100--120 & 1 & 11 & 2045 & 0 & 2 & 1 & incrementa \\
120--140 & 1 & 11 & 2045 & 0 & 3 & 2 & incrementa \\
140--160 & 1 & 11 & 2045 & 0 & 4 & 3 & incrementa \\
160--180 & 1 & 00 & 2045 & 0 & 3 & 3 & mantém \\
180--200 & 1 & 00 & 2045 & 0 & 3 & 3 & mantém \\
200--220 & 1 & 00 & 2045 & 0 & 3 & 3 & mantém \\
\hline
\end{longtable}}
\vspace{-0.9cm}
\fonte{Autoria própria. Saídas do fim de cada intervalo, conforme o modelo de referência e a simulação do Quartus II.}

{\scriptsize\setlength{\tabcolsep}{3pt}
\begin{longtable}{rccccccp{4.2cm}}
\caption{Casos de teste de \texttt{prog\_cnt\_carga.vwf}: cargas e next\_pc\_out combinacional.}
\label{tab:caso-prog-cnt-carga}\\
\hline $t$ (ns) & \texttt{nrst} & \texttt{pc\_ctrl} & \texttt{new\_pc\_in} & \texttt{from\_stack} & \texttt{next\_pc\_out} & \texttt{pc\_out} & Observação \\ \hline \endfirsthead
\hline $t$ (ns) & \texttt{nrst} & \texttt{pc\_ctrl} & \texttt{new\_pc\_in} & \texttt{from\_stack} & \texttt{next\_pc\_out} & \texttt{pc\_out} & Observação \\ \hline \endhead
0--20 & 0 & 00 & 1000 & 2000 & 0 & 0 &  \\
20--40 & 1 & 10 & 1000 & 2000 & 1000 & 1000 & pc\_ctrl = 10: carga de new\_pc\_in = 1000 (from\_stack = 2000 é ignorado) \\
40--60 & 1 & 01 & 1000 & 2000 & 2000 & 2000 & pc\_ctrl = 01: carga de from\_stack = 2000 \\
60--65 & 1 & 00 & 1000 & 2000 & 2000 & 2000 & pc\_ctrl = 00: mantém \\
65--80 & 1 & 00 & 1 & 2000 & 2000 & 2000 & new\_pc\_in = 1: sem efeito (pc\_ctrl = 00) \\
80--105 & 1 & 10 & 1 & 2000 & 1 & 1 & pc\_ctrl = 10: carga de new\_pc\_in = 1 \\
105--120 & 1 & 01 & 1 & 2000 & 2000 & 2000 & pc\_ctrl = 01 antes da borda: next\_pc\_out muda na hora para from\_stack \\
120--140 & 1 & 11 & 1 & 2000 & 2002 & 2001 & pc\_ctrl = 11: next\_pc\_out = pc + 1 \\
140--145 & 1 & 11 & 1 & 2000 & 2002 & 2001 & pc\_ctrl = 00 \\
145--165 & 1 & 00 & 1 & 2000 & 2001 & 2001 & pc\_ctrl = 00 e pc = 2001 (hold) \\
165--185 & 1 & 00 & 1 & 1365 & 2001 & 2001 & from\_stack muda para 1365: next\_pc\_out só reage com pc\_ctrl = 01 \\
185--205 & 1 & 01 & 1 & 1365 & 1365 & 1365 & pc\_ctrl = 01 \\
205--225 & 1 & 10 & 682 & 1365 & 682 & 682 & new\_pc\_in = 682, pc\_ctrl = 10 \\
225--245 & 1 & 11 & 682 & 1365 & 684 & 683 & pc\_ctrl = 11 \\
245--265 & 1 & 00 & 682 & 1365 & 683 & 683 & pc\_ctrl = 00 \\
265--285 & 1 & 01 & 682 & 0 & 0 & 0 & carga de 0 pela pilha \\
285--305 & 1 & 10 & 2047 & 0 & 2047 & 2047 & carga de 2047 por new\_pc\_in \\
305--325 & 1 & 11 & 2047 & 0 & 1 & 0 & incrementa: volta a 0 \\
325--360 & 1 & 00 & 2047 & 0 & 0 & 0 & mantém \\
\hline
\end{longtable}}
\vspace{-0.9cm}
\fonte{Autoria própria. Saídas do fim de cada intervalo, conforme o modelo de referência e a simulação do Quartus II.}

{\scriptsize\setlength{\tabcolsep}{3pt}
\begin{longtable}{rccccccp{4.2cm}}
\caption{Casos de teste de \texttt{prog\_cnt\_reset.vwf}: reset assíncrono.}
\label{tab:caso-prog-cnt-reset}\\
\hline $t$ (ns) & \texttt{nrst} & \texttt{pc\_ctrl} & \texttt{new\_pc\_in} & \texttt{from\_stack} & \texttt{next\_pc\_out} & \texttt{pc\_out} & Observação \\ \hline \endfirsthead
\hline $t$ (ns) & \texttt{nrst} & \texttt{pc\_ctrl} & \texttt{new\_pc\_in} & \texttt{from\_stack} & \texttt{next\_pc\_out} & \texttt{pc\_out} & Observação \\ \hline \endhead
0--20 & 0 & 00 & 1234 & 0 & 0 & 0 &  \\
20--40 & 1 & 10 & 1234 & 0 & 1234 & 1234 & carrega 1234 \\
40--60 & 1 & 11 & 1234 & 0 & 1236 & 1235 & incrementa \\
60--65 & 1 & 11 & 1234 & 0 & 1236 & 1235 & incrementa \\
65--80 & 0 & 11 & 1234 & 0 & 1 & 0 & reset assíncrono, fora de qualquer borda \\
80--100 & 0 & 11 & 1234 & 0 & 1 & 0 & ainda em reset; incremento sobre a borda: não conta \\
100--123 & 0 & 10 & 1234 & 0 & 1234 & 0 & ainda em reset; carga sobre a borda: não carrega \\
123--140 & 1 & 10 & 1234 & 0 & 1234 & 1234 & solta o reset \\
140--160 & 1 & 11 & 1234 & 0 & 1236 & 1235 & incrementa (a borda desse ciclo já conta) \\
160--192 & 1 & 11 & 1234 & 0 & 1238 & 1237 & incrementa \\
192--195 & 0 & 11 & 1234 & 0 & 1 & 0 & pulso de reset de 3 ns, sem borda \\
195--200 & 1 & 11 & 1234 & 0 & 1 & 0 & fim do pulso \\
200--220 & 1 & 11 & 1234 & 0 & 2 & 1 & incrementa a partir de 0 \\
220--240 & 1 & 00 & 1234 & 0 & 1 & 1 & mantém \\
\hline
\end{longtable}}
\vspace{-0.9cm}
\fonte{Autoria própria. Saídas do fim de cada intervalo, conforme o modelo de referência e a simulação do Quartus II.}

